
\section{Eighteenth Sunday after Trinity: What Is Impossible for Men is Possible for God}

\begin{quote}

Mark 10:17–27. And as He had gone forth onto the road, one came running and fell on his knees before Him and asked Him: Good Master! What shall I do, that I may inherit eternal life? But Jesus said to him: Why do you call Me good? None is good save One, that is, God. You know the commandments: You shall not commit adultery, you shall not murder, you shall not steal, you shall not bear false witness, you shall not defraud, honor your father and your mother! But he answered and said to Him: Master! All these have I kept from my youth. But Jesus looked upon him and loved him, and said to him: One thing you lack: go away, sell what you have, and give it to the poor, and you shall have a treasure in heaven; then come, follow Me, and take up the cross! But he was troubled at that saying and went away sorrowful; for he had great possessions. And Jesus looked around and said to His disciples: How difficult it shall be for those who have riches to enter into the kingdom of God! But the disciples were astonished at His words. But Jesus answered again and said to them: Children! How difficult it is for those who trust in their riches to enter into the kingdom of God! It is easier for a camel to go through a needle’s eye than for a rich man to enter into the kingdom of God. But they were exceedingly astonished, and said among themselves: Who then can be saved? But Jesus looked upon them and said: For men it is impossible, but not for God; for all things are possible for God.

\end{quote}

\bigskip

When God does the impossible, it is not in order to reveal His omnipotence, — for that He does also through the regular order of nature (Rom. 1:20), — but in order to save men, and to reveal His love.

Sarah laughed when God promised Abraham offspring by her; Abraham himself, on the other hand, looked neither to his old age nor to Sarah’s dead womb, but in childlike faith held fast to the Promise, that in his seed all the families of the earth should be blessed; and as Abraham believed, so it came to pass; what was impossible for men, that God did; He gave Sarah a son, and thus initiated the fulfillment of the great Promise: “The woman’s seed shall crush the serpent’s head.”

For the sake of this same cause God did what, so to speak, is yet more impossible for men: He caused a Virgin to bear a Child without knowing man. “The Holy Spirit shall come upon you, and the power of the Most High shall overshadow you; therefore also that Holy One who is born of you shall be called the Son of God.” That was the angel’s message to Mary, and according to this word God’s own Son was born of a woman, under the Law, subjected to all the conditions of corruptibility, of the same flesh and blood as you and I, and yet of another order, a Man without sin, who could die for us, overcome death, and deliver us from sin and death and the devil, in order that all God’s promise of enmity should be fulfilled in the woman’s seed, the Son of God and of Mary.

God’s love does not stop here; from this point of departure it so to speak only then begins truly to reveal its power, which, as the Word says, is stronger than death. There is one thing which, if one may so express it, is a thousandfold more impossible than anything else: that a sinful, corruptible human child should be able to receive a new nature, undergo a new birth, and thereby become a child of God, a partaker of the divine nature, holy and incorruptible in the midst of a world stained by sin and in a corruptible body. With good reason Nicodemus could ask: “How can a man be born when he is old? Can he enter a second time into his mother’s womb and be born?” And would that very many might ask the question with the same seriousness as Nicodemus, and not give it up until they had received an answer from God, until they had been born “not of blood, nor of the will of the flesh, nor of the will of man, but of God,” and had truly learned by experience that “what is impossible for men is possible for God.”

For “one thing you still lack,” the Lord says to the rich young man, one thing you lack; it is this one thing that makes it impossible for a man to inherit eternal life. One thing you trust in, be it never so slight, never so corruptible, — that sets an impassable wall between you and your God, between your soul and eternal life.

Friend, have you found in yourself this one thing which is lacking and stands in the way of your salvation?

“Good Master, what shall I do, that I may inherit eternal life?” cries out the rich young man.

He had striven so honestly and earnestly; the commandments he had kept from his youth, — how many can say the same? — and yet he was no nearer. Now he came, surely after much inward struggle, — for he was a ruler (Luke 18:18), — to the despised Nazarene, who was even then on the way to His suffering, and asked His counsel. That he also asked in all sincerity seems to follow from what is said of Jesus: “He looked upon him and loved him.” And yet, when he received this answer: “One thing you still lack; sell what you have, and give it to the poor; then come and follow Me and take up the cross!” then he went away sorrowful, and it is added: “For he had great possessions.”

This man lacked neither earnest striving nor a sincere desire to be saved; he wanted so gladly to be saved, but he could not; there was one thing he could not give up; therefore he could not follow Jesus and take up the cross either.

Alas, how many are there not still to this very day for whom the cross thus becomes an offence, and who in despondency turn away from Jesus, if not to the world’s frivolity and vices, then at least to its outward righteousness! A little stone in their path became a mountain and shut them out from God.

And God’s children are daily exposed to the same danger: to forget and cast away the cross in order at once to seize the crown. A seemingly little thing comes into their way, clings like a thorn to their hearts, and pierces them with death’s sting.

Therefore Jesus looked around upon His disciples when the young man went away. He knew what thoughts were stirring in their hearts. The hope of an outward glorious Messianic kingdom still had power over their souls. Is it not possible, then, to be rich in the kingdom of God? Is not the kingdom of God a kingdom of riches and power? Thus they thought, more or less dimly consciously, and were therefore also astonished when Jesus told them that it is exceedingly difficult for a rich man to enter into the kingdom of God. And when He both explained and sharpened this word to mean that rich in this sense is the one who trusts in his riches, and that it is just as impossible for such a one to enter into the kingdom of God as for a camel to go through a needle’s eye, — then, it is written, they became even more exceedingly astonished, and asked one another: “Who then can be saved?”

Staggering disciples! As Peter said, they had left all to follow Jesus (v. 28), in order thereby to attain the glory of the Messianic kingdom according to the carnal Jewish conception they had of it.

Now it was as though everything had been torn from their hands. Jesus was to suffer the reproach of the cross. The disciples would indeed also have earthly things in the world, but all under persecution, and besides this “many who were first shall be last” (vv. 30, 31). Even for the Lord’s disciples, then, only one thing here below is certain: the cross. There, it is eternal life, that is true, but only for the one who remains faithful unto the end. And who is sufficient for that? When the one thing, sometimes of a purely earthly sort, such as possessions, honor, influence, family, fleshly desires and thoughts, — just as often of a spiritual sort, such as anger, envy, hatred, sloth, esteem as a good Christian, yes, even inward vanity over one’s Christian strength, — can hinder and bind and capture and destroy the soul of a child of God, oh, how must we not, like the disciples, be filled with fear and ask: “who then can be saved?” and anew humbly flee to Golgotha and surrender ourselves wholly and altogether into His hands, — blessed be His name forever! — for whom that is possible which is impossible for us poor men: to make our soul righteous before God!

But if the righteous is scarcely saved, where shall the sinner and the ungodly appear?

Brother and Sister! Have you, with the rich young man, turned away from Jesus because His cross frightened you, and because He required of you the one thing which it was impossible for you to give up? Where then are you going? Think seriously and simply about this question: Where does it lead? And when you are finished, when with horror you have looked down into the abyss where your existence will be transformed into eternal torment, and in anguish begin to look back toward the Father in heaven and see an impassable wall between yourself and the Father’s heart, oh, then look to Him from whose thorn-crowned brow the drops of blood are running; see, He goes out into the wilderness precisely to seek you, He does not wait until you come to Him, He seeks you, He stands at your heart’s door, He knocks, He will come in and whisper this one word into your despairing soul:

“What is impossible for men is possible for God.”

Look to Him, the Savior, born of Abraham’s seed, of a Virgin, dead, risen, ascended into heaven! Can He not save?

